% BANKS-SOURCE: pdf=91 print=81 kind=section
\subsection{Spontaneously broken symmetries}

% BANKS-SOURCE: pdf=91 print=81 kind=prose
The statement that the free massless spin-$1$ Belinfante tensor is traceless
is four-dimensional, $D=4$; the massless spin-$1/2$ statement extends to
general $D$. This is not the case for spin zero, but we can write a new tensor
\begin{equation*}
  T_I^{\mu\nu}=T_B^{\mu\nu}
  -\xi_D(\partial^\mu\partial^\nu-\eta^{\mu\nu}\partial^2)\phi^2,
  \qquad \xi_D=\frac{D-2}{4(D-1)}.
\end{equation*}
which is traceless on shell.  In $D=4$, $\xi_D=1/6$.  Note that this improved
stress tensor is not invariant under
translation of the field $\phi\to\phi+c$, whereas the Belinfante tensor is.\footnote[4]{The diligent reader should take all the unproven statements in this
section as additional exercises.} There is an interesting connection between
the lack of conformal invariance of the spin-zero Belinfante tensor and the
theory of Nambu--Goldstone bosons (NGBs) that we will investigate in the next
section. The field translation invariance of the massless free scalar is the
simplest example of a spontaneously broken symmetry. We will learn that every
independent one-parameter group of symmetries that does not leave the vacuum
state invariant gives rise to a massless spin-zero particle, the NGB. The theory
of NGBs contains a dimensionful constant, $f$, defined by
\begin{equation*}
  \langle0|J^\mu(0)|p\rangle
  =\frac{fp^\mu}{\sqrt{(2\pi)^{D-1}2\omega_{\mathbf p}}}.
\end{equation*}
In $D$ space-time dimensions currents have mass dimension $D-1$ while states
have mass dimension $(1-D)/2$. Thus $f$ has mass dimension $(D-2)/2$. The theory
of NGBs is not conformally invariant, except perhaps for $D=2$.\footnote[5]{In
fact, the Coleman--Mermin--Wagner theorem tells us that spontaneous breakdown
of compact symmetry groups does not occur in two dimensions. In particular, the
abelian NGB theory is conformally invariant, and does not have spontaneous
symmetry breakdown.} Correspondingly, the traceless stress tensor for the
massless field explicitly breaks the field translation symmetry.

\BanksImplicitHook{I-CH07-003}

To summarize, the W--T identities for Poincare symmetry are encoded in the
conservation and symmetry of the Belinfante stress tensor. Conformally invariant
theories satisfy the additional property that the stress tensor is traceless.
When we study renormalization in a later section, we will learn that classical
conformal invariance usually fails in the quantum theory. Conformal quantum
field theories (CFTs) are few and far between, but they provide us with the
proper definition of all quantum field theories in existence.

% BANKS-SOURCE: pdf=91 print=81 kind=prose
The predictions of symmetries are powerful and have broad applications.
Surprisingly, even more wonderful results appear when the symmetry is broken
spontaneously. We use this phrase to describe what happens when, despite an
exact symmetry of the Lagrangian, the Green functions of a theory do not obey
the symmetry predictions we have just discussed. This is a phenomenon special
to quantum field theory, resulting from the infinite volume of space. Of course,
in the real world we do not know that the volume is infinite. Nonetheless, the
predictions of spontaneously broken symmetry are valid with high precision even
for finite volume as long as the dynamics is governed by local field theory,
perhaps with an ultraviolet cut-off, and the volume is large in cut-off units.

To see how spontaneously broken symmetry can occur, we write the Noether
formula for variable $\epsilon$ and take the limit of constant $\epsilon$ more
carefully. Thus, by differentiating $k$ times with respect to $J$, we obtain
\begin{equation*}
  \partial^x_\mu\langle J^\mu(x)\phi_{i_1}(x_1)\cdots\phi_{i_k}(x_k)\rangle
  =\ii\sum_n\delta^D(x-x_n)\langle\phi_{i_1}(x_1)\cdots
  T_{j i_n}\phi_j(x_n)\cdots\phi_{i_k}(x_k)\rangle.
\end{equation*}
This is called the Ward--Takahashi identity.

On integrating this with respect to $x$ and dropping the integral of a total
derivative, we obtain the invariance of the Green functions. The derivation can
fail if it is illegitimate to drop the integrated total derivative. If we take
all of the $x_i$ to the same point $y$ (in a way we will understand better when
we discuss renormalization -- this is called an operator product expansion),
invariance will fail if and only if
\begin{equation*}
  \int\dd^D x\,\partial_\mu\langle J^\mu(x)\mathcal O_A(y)\rangle\ne0,
\end{equation*}
for some (possibly composite) operator $\mathcal O_A$, which transforms in a
non-trivial representation of the symmetry group. In Fourier space this is the
statement that
\begin{equation*}
  p_\mu\Gamma_A^\mu(p)=p^2\Gamma_A(p^2)\ne0,
\end{equation*}
when $p\to0$. Here we have used Lorentz invariance to write the Fourier
transform of the Green function, $\Gamma_A^\mu(p)=p^\mu\Gamma_A(p^2)$. We
conclude that the two-point function has a pole, so the theory contains a
massless particle that can be created from the vacuum both by the current and by
the field $\mathcal O_A$. This is the celebrated Nambu--Goldstone boson
\cite{banks-ref-36,banks-ref-37,banks-ref-38,banks-ref-39,banks-ref-40}. It carries spin zero, because the right-hand side of the equation can
be non-vanishing only if $\langle\mathcal O_A\rangle\ne0$. If Lorentz
invariance is not itself spontaneously broken (which does not happen in most
field theories), then $\mathcal O_A$ is a scalar field. Notice that this also
assumes that the symmetry is an internal symmetry, which does not change the spin
of fields. In space-time dimensions higher than 2, the only exception to this
rule is supersymmetry. In that case the Nambu--Goldstone (NG) particle is a
fermion, called the Goldstino.

% BANKS-SOURCE: pdf=92 print=82 kind=prose
We can obtain a more physical understanding of this result, and one that is
useful in systems that are not Lorentz-invariant, by thinking about the
implications of the non-zero vacuum expectation value (VEV) in the operator
formalism. In operator language, a one-parameter group of symmetries is
implemented by a one-parameter group of unitary operators
$U(\omega)=e^{\ii\omega Q}$, where $Q$ is Hermitian. A symmetry-violating VEV
is possible only if the vacuum is not invariant under the symmetry
transformation, which means that $Q|0\rangle\ne0$. On the other hand, since
$Q$ commutes with the Hamiltonian, the states $U(\omega)|0\rangle$ all have
zero energy.\footnote[6]{We are ignoring mathematical rigor here. In fact the
operators $U(\omega)$ are not even valid operators in the Hilbert space of our
system. But the argument is valid nonetheless.} Now consider instead the
operator
\begin{equation*}
  Q_V=\int\dd^{D-1}\mathbf x\,\omega(\mathbf x)J^0(x),
\end{equation*}
where $\omega(\mathbf x)$ is a smooth function which vanishes outside of a volume $V$
and is equal to a constant over most of the inside of this volume. We will
consider taking the volume to infinity and letting the function $\omega(\mathbf x)$
become more and more slowly varying. The commutator of the Hamiltonian with
this operator is (we use current conservation)
\begin{equation*}
  [H,Q_V]\propto\int\dd^{D-1}\mathbf x\,\boldsymbol\nabla\omega\cdot\mathbf J,
\end{equation*}
which can be made arbitrarily small as we take $V\to\infty$ and
$\boldsymbol\nabla\omega\to0$. Thus we find, as a consequence of spontaneously
broken continuous symmetry (a symmetry that does not preserve the ground state),
a spectrum of localized excitations of arbitrarily low energy. These are the NG
bosons. In a relativistic theory they must be spin-zero massless particles. The
phenomenon of spontaneous symmetry breaking occurs in non-relativistic
condensed-matter systems. In general these systems may not have any space-time
symmetries, so all we can conclude is that the NG bosons have a dispersion
relation satisfying $E_{\mathbf p}\to0$ as $|\mathbf p|\to0$. (Note that the energy gap in a
non-relativistic system has nothing to do with particle masses.)

The picture to keep in mind for understanding what is going on in a theory with
spontaneous symmetry breaking is shown in Figures~\ref{fig:7.1} and~\ref{fig:7.2}.
The first of these shows a ground state with spontaneously broken symmetries.
The rotational symmetry of the arrows is broken because they all want to point
in the same direction. Figure~\ref{fig:7.2} shows a Nambu--Goldstone excitation
of such a ground state: a long-wavelength wave in the directions of the arrows.

% BANKS-SOURCE: pdf=93 print=83 kind=figure id=7.1
\begin{figure}[htbp]
  \centering
  % BANKS-SOURCE: pdf=93 print=83 kind=figure id=7.1
\resizebox{10.5cm}{!}{%
\begin{tikzpicture}[x=1cm,y=1cm,
    BanksSpin/.style={line width=1pt,-{Stealth[length=8.8pt,width=5.6pt]}}]
  % 24 arrows on a single row, all pointing straight up
  \foreach \k in {0,...,23} {
    \draw[BanksSpin] ({0.415*\k},0) -- ({0.415*\k},2.00);
  }
  % ellipsis dots: half-spacing steps beyond each end of the row
  \foreach \k in {-1.5,-1,-0.5,23.5,24,24.5} {
    \fill ({0.415*\k},0.051) circle (1.45pt);
  }
\end{tikzpicture}%
}

  \caption{A state with spontaneously broken symmetry.}
  \label{fig:7.1}
\end{figure}

% BANKS-SOURCE: pdf=93 print=83 kind=figure id=7.2
\begin{figure}[htbp]
  \centering
  % BANKS-SOURCE: pdf=93 print=83 kind=figure id=7.2
\resizebox{10.5cm}{!}{%
\begin{tikzpicture}[x=1cm,y=1cm,
    BanksSpin/.style={line width=1pt,-{Stealth[length=8.8pt,width=5.6pt]}}]
  % Same 24-site row as Figure 7.1, but the arrow directions sweep smoothly
  % down from vertical on the left to about 24 degrees just left of centre and
  % then back up to vertical on the right (a long-wavelength wave).
  \foreach \k/\ang in {%
    0/90, 1/81, 2/73, 3/59.5, 4/50.5, 5/42.5, 6/34.5, 7/29,
    8/26.5, 9/24, 10/26.5, 11/26.5, 12/28, 13/31, 14/34.5, 15/40,
    16/46, 17/53.5, 18/57.5, 19/65, 20/70.5, 21/77, 22/83.5, 23/90} {
    \draw[BanksSpin] ({0.415*\k},0) -- ++(\ang:1.93);
  }
  % ellipsis dots: half-spacing steps beyond each end of the row
  \foreach \k in {-1.5,-1,-0.5,23.5,24,24.5} {
    \fill ({0.415*\k},0.051) circle (1.45pt);
  }
\end{tikzpicture}%
}

  \caption{A Nambu--Goldstone excitation.}
  \label{fig:7.2}
\end{figure}

A final note: it is not necessary for an elementary field, which appears in the
Lagrangian, to have a VEV. If there is any Green function of elementary fields
that

% BANKS-SOURCE: pdf=94 print=84 kind=prose
violates a symmetry, we can consider the limit where all points in this function
are taken close to the same point $x$. As we will discuss later, in this limit the
product of elementary fields has an expansion (operator product expansion or
OPE) in terms of functions of coordinate differences, multiplied by local
composite operators evaluated at $x$. For example, in free field theory
\begin{equation*}
  \phi(x)\phi(0)\sim\frac{1}{(x-y)^2}\mathcal A
  +\mathcal B:\!\phi^2(0)\!:+o(x-y).
\end{equation*}
If a symmetry-violating Green function is non-vanishing then some local
composite operator $\mathcal O_A$ will have a non-zero VEV, and we can replay
Goldstone's theorem with $S_A\to\mathcal O_A$. So, like the LSZ formula,
Goldstone's theorem applies to composites of Lagrangian fields, and is more
general than the Lagrangian formalism.
